\documentclass[10pt,a4paper]{amsart}
\usepackage[margin=19mm]{geometry}
\usepackage{amsmath,amssymb,mathtools}
\usepackage[scr=rsfs,cal=euler]{mathalfa}
\usepackage{xfrac}
\usepackage[unicode,hidelinks,bookmarks=false]{hyperref}
\newcommand{\ie}{i.e.\ }
\newcommand{\cf}{cf.\ }
\newcommand{\resp}{respectively\ }
\newcommand{\Subsection}{\subsection}
\providecommand{\nobreakdash}{\nobreak\hbox{-}}
\setlength{\emergencystretch}{2em}
\multlinegap0pt
\usepackage{bm}
\usepackage{bbm}
\usepackage{dsfont}
\usepackage{xspace}
\usepackage{cancel}
\usepackage{mathtools}
\usepackage{cleveref}
\usepackage{amsthm}
\makeatletter
\@ifundefined{thm}{\newtheorem{thm}{Thm}}{}
\@ifundefined{cor}{\newtheorem{cor}[thm]{Corollary}}{}
\@ifundefined{lem}{\newtheorem{lem}[thm]{Lemma}}{}
\makeatother
\DeclareMathOperator{\C}{\mathbb{C}}
\DeclareMathOperator{\Imagine}{Im}
\DeclareMathOperator{\R}{\mathbb{R}}
\DeclareMathOperator{\SO}{SO}
\title[A generalized Frankel conjecture via the Yang-Mills flow]{A generalized Frankel conjecture via the Yang-Mills flow}
\author{Jiangtao Li}
\date{Source: arXiv:2511.04003v1 (CC BY 4.0)}
\begin{document}
\maketitle

\noindent\emph{Source and licence.} A generalized Frankel conjecture via the Yang-Mills flow. Version arXiv:2511.04003v1 (CC BY 4.0), distributed under the Creative Commons Attribution 4.0 International licence, as recorded in the arXiv licence metadata for this exact version and shown on the abstract page https://arxiv.org/abs/2511.04003v1. Full rights notice: licenses/arxiv-2511.04003-CC-BY-4.0.txt.\par\smallskip

\noindent\textit{Editorial scope.} Counted unit: the Cor whose source label is \texttt{None} in the article (source0.tex lines 256--258, source label \texttt{None}), delivered together with the article's own complete original proof of it (source0.tex lines 259--293, one \texttt{proof} environment of 35 lines). No mathematics is written, repaired or re-derived: statement and proof are the article's own text line for line. Required context retained uncounted: lem:BisectionalCurvature (lines 233--246); eq:BisectionalCurvature (lines 243--245). Every definition, display and label of those lines is retained unchanged. Omitted from this package and not redistributed: 1--232, 246--255, 294--498. Nothing inside a retained excerpt is omitted or altered: every retained body is delivered exactly as the article's lines are written, so no omission receipt of any kind accompanies this package. Citations: the retained excerpts carry no citation command at all, so no bibliography entry of the article is transcribed and no reference is redistributed. Reference adaptation: none was needed and none was made; every reference inside the delivered unit resolves inside it, so no unresolved reference or citation is produced. Printed slips kept literal: no printed word, symbol, hypothesis or number of the counted statement or of its proof was altered. Layout and class: the article's own class is \texttt{amsart} with packages including amsmath, amssymb, comment, color, xcolor, soul, geometry, hyperref, cleveref, quiver, esint, nicematrix; the wrapper is the lane's pinned \texttt{amsart} editorial wrapper at 10pt on A4, which restates the theorem-like environments the retained text needs and adds the article's own macro definitions that the retained text needs (\texttt{\textbackslash C}, \texttt{\textbackslash Imagine}, \texttt{\textbackslash R}, \texttt{\textbackslash SO}), with extra wrapper packages (bm, bbm, dsfont, xspace, cancel, mathtools, cleveref). The delivered numbering is the unit's own. Assets: no retained span contains a figure environment, a graphics-inclusion command, a PDF-page inclusion, a table or a TikZ picture; no image file is copied into this package, the package declares no asset dependency and no image file is redistributed with it. Licence: version arXiv:2511.04003v1 of this article is distributed under the Creative Commons Attribution 4.0 International licence, as recorded in the arXiv licence metadata for this exact version and shown on the abstract page https://arxiv.org/abs/2511.04003v1. A full rights notice is delivered at licenses/arxiv-2511.04003-CC-BY-4.0.txt. Characters: every retained excerpt is pure ASCII, so the delivered engine, XeLaTeX with its default Latin Modern text font, reports no missing character.
\begin{lem}\label{lem:BisectionalCurvature}
    Let 
    \[U=\begin{pmatrix}
        a_1&\cdots&a_n\\
        -\sqrt{-1}a_1&\cdots&-\sqrt{-1}a_n
    \end{pmatrix},\quad V=\begin{pmatrix}
        b_1&\cdots&b_n\\
        -\sqrt{-1}b_1&\cdots&-\sqrt{-1}b_n
    \end{pmatrix}\]
    be holomorphic tangent vectors. Then 
    \begin{equation}\label{eq:BisectionalCurvature}
        R(U,\Bar{U},V,\Bar{V})=4\left(\sum_{i=1}^n|a_i|^2\right)\left(\sum_{i=1}^n|b_i|^2\right)-16\sum_{i<j}\Imagine(\Bar{a_i}a_j)\Imagine(b_i\Bar{b_j}).
    \end{equation}
\end{lem}

\begin{cor}
    The Hermitian symmetric space $Q^n$ has nonnegative bisectional curvature. More precisely, for any holomorphic tangent vector $U\in T^{1,0}_xQ^n$, $R(U,\Bar{U}):T_x^{1,0}Q^n\to T_x^{1,0}Q^n$ has nonnegative eigenvalues and at most one of them is zero. In particular, $Q^n$ has $2$-positive bisectional curvature.
\end{cor}

\begin{proof}
    Observe that the curvature is invariant under the adjoint action of the isotropy group $\SO(n,\R)\times\SO(2,\R)$. Therefore, we have 
    \begin{equation}\label{eq:InvarianceOfCurvature}
    R(BUA,B\Bar{U}A,BVA,B\Bar{V}A)=R(U,\Bar{U},V,\Bar{V}), \quad \forall A\in\SO(n,\R), B\in\SO(2,\R).
    \end{equation}
    Suppose that $U,V$ are nonzero holomorphic tangent vectors as in Definition \ref{lem:BisectionalCurvature}. By \ref{eq:InvarianceOfCurvature}, we can assume, without loss of generality, that $a_i$'s are all zero except for $a_1$ and $a_2$. Next we split into two cases:
    \begin{enumerate}
        \item If one of $a_1$ and $a_2$ is zero, then it follows from \Cref{eq:BisectionalCurvature} that 
        \[R(U,\Bar{U},V,\Bar{V})=4\left(\sum_{i=1}^n|a_i|^2\right)\left(\sum_{i=1}^n|b_i|^2\right)> 0.\]
        \item If both $a_1$ and $a_2$ are nonzero, then 
        \begin{align*}
            &R(U,\Bar{U},V,\Bar{V})\\
            =&4\left(\sum_{i=1}^n|a_i|^2\right)\left(\sum_{i=1}^n|b_i|^2\right)-8\sum_{i\ne j}^n\Imagine(\Bar{a_i}a_j)\Imagine(b_i\Bar{b_j})\\
            =&4\left(\sum_{i=1}^n|a_i|^2\right)\left(\sum_{i=1}^n|b_i|^2\right)-16\Imagine(\Bar{a_1}a_2)\Imagine(b_1\Bar{b_2})\\
            \ge&4\left(\sum_{i=1}^n|a_i|^2|b_i|^2+\sum_{i\ne j}|a_i|^2|b_j|^2\right)-16|a_1||a_2||b_1||b_2|\\
            =&4\left(\sum_{i=1}^2|a_i|^2|b_i|^2-2|a_1||b_1||a_2||b_2|\right)+4\left(\sum_{j\ne 1}|a_1|^2|b_j|^2+\sum_{j\ne 2}|a_2|^2|b_j|^2-2|a_1||b_1||a_2||b_2|\right)\\
            =&4(|a_1||b_1|-|a_2||b_2|)^2+(|a_1||b_2|-|a_2||b_1|)^2+4\left(\sum_{j\ne 1,2}|b_j|^2\right)(|a_1|^2+|a_2|^2)\\
            \ge&0.
        \end{align*}
        The equality holds if and only if 
        \begin{equation*}
            U=
            \begin{pmatrix}
                a&\sqrt{-1}a&\cdots&0\\
                -\sqrt{-1}a&a&\cdots&0
            \end{pmatrix},
            \quad
            V=\begin{pmatrix}
                b&-\sqrt{-1}b&\cdots&0\\
                -\sqrt{-1}b&-b&\cdots&0
            \end{pmatrix}, \quad \forall a,b\in\C^\times.
        \end{equation*}
    \end{enumerate}
    It follows from the above argument that $R$ has nonnegative and $2$-positive bisectional curvature.
\end{proof}

\end{document}
